\section{Background and Motivation}

A private investor holding an S\&P 500 asset faces a similar decision
each trading day: buy more, sell part of the position, or do nothing. On
most days, the choice may have little effect on the overall outcome.
During a heavy market move, however, the decision can have a much larger
impact.

March 2025 is one of such periods studied in this dissertation. Under
the experimental setup used later in this study, it was the worst month
for a simple buy-and-hold strategy, which lost \$398.41 from an initial
\$10,000 investment, approximately 4\% of the starting capital. An investor who reduced their position before the decline could
have avoided part of this loss. However, reducing positions too often
would also mean missing gains during the other months, which produced an
average monthly gain of more than \$200.

The problem is therefore not just when to invest, or deciding when the
market will rise or fall, but how to make a series of decisions about an
asset's position as market conditions change, as each decision made
affects the position long-term and the consequences of that action will
only become clear several days or weeks later.

This makes trading a sequential decision-making problem, which is one
reason reinforcement learning (RL)~\cite{sutton2018} is appropriate for the task. RL
allows an agent to learn by interacting with an environment, observing
the results of its actions, and receiving rewards over time based on
those results. In essence, it learns a policy/strategy that makes better
decisions to receive more rewards.

Trading has several characteristics that make it suitable for this
machine learning approach. The agent observes information about the
market and its portfolio, selects from a set of possible actions such as
buying, selling, or holding, and receives a measurable reward based on
the change in portfolio wealth. As a result, reinforcement learning has
been increasingly applied to financial trading~\cite{moody2001, deng2017, fischer2018}.

However, another important problem arises. The Stock/Equity markets
generally rise over the years rather than fall. Meaning that, in a
rising market, a strategy that simply remains invested without any
pulling out can generate a strong return without using any market
information. A learned agent that produces a positive return therefore
has not necessarily demonstrated, by itself, that it has learned a
useful trading strategy in such markets.

Two questions become relevant. First, can the learned agent outperform
the simple buy-and-hold strategy with which it is being compared?
Second, does the agent actually use the information available to it when
deciding what action to take? These questions are important when
evaluating on data not used during training.

This dissertation therefore focuses on both the financial performance of
reinforcement learning agents and whether their decisions show evidence
of using the information provided to them.

\section{Problem Statement}

The main problem this dissertation addresses is whether reinforcement
learning algorithms can learn a trading policy for a portfolio ---
containing only one asset in this case --- that beats a simple benchmark
strategy on unseen data.

However, financial performance alone is not enough to determine that an
agent has actually learned a useful trading strategy. For example, an
agent that buys whenever it can, would perform well in a generally
rising market without using any of the market information provided to
it. A positive return from such a policy could be incorrectly
interpreted as proof that the agent learned when to buy, sell, or remain
out of the market, whereas that was not the case.

Evaluating a learned trading agent therefore requires a second
perspective alongside financial performance: \textbf{behaviour}. The
study must analyse whether the agent\textquotesingle s actions are a
result of changes in the observed market and portfolio state. This helps
to differentiate a policy that genuinely responds to the information it
receives from one that simply follows a fixed trading pattern.

This dissertation therefore addresses two interrelated problems: whether
reinforcement learning can produce better financial performance than
simple benchmarks, and whether that performance stems from a learned
policy that actually uses the information available to it.

\section{Research Questions and Objectives}

The aim of this dissertation is to evaluate three reinforcement learning
algorithms --- REINFORCE~\cite{williams1992}, Deep Q-learning (DQN)~\cite{mnih2015}, and Proximal Policy
Optimization (PPO)~\cite{schulman2017ppo} --- and determine whether they can learn profitable
and genuinely intelligent trading strategies for the SPY exchange-traded
fund.

The study is structured around four research questions:

\begin{enumerate}
\def\labelenumi{\arabic{enumi}.}
\item
  \textbf{Performance:} Do the learned agents outperform simple
  benchmarks on unseen data after transaction costs are taken into
  account?
\item
  \textbf{Behaviour:} Do the learned agents use the market and portfolio
  information provided in their state representation when making trading
  decisions?
\item
  \textbf{State representation:} Does adding a drawdown feature to the
  main nine-feature state change the behaviour learned or the financial
  performance?
\item
  \textbf{Reward:} Does a risk-aware reward improve the balance between
  return and drawdown, rather than simply encouraging the agent to trade
  less?
\end{enumerate}

These research questions are touched upon through seven measurable
objectives, each traceable to a specific part of the dissertation

\textbf{O1.} Build a reproducible daily trading environment for SPY
using monthly episodes, transaction fees, and action masking, and verify
the environment (Chapter 4).

\textbf{O2.} Implement REINFORCE, DQN, and PPO and verify their
implementations before running the main trading experiments (Chapters 3
and 4).

\textbf{O3.} Calibrate the fixed trade size on validation data so the
agents' exposure to the market is not limited, providing a fair
comparison to the buy-and-hold benchmark (Section 5.3).

\textbf{O4.} Apply behavioural analysis to determine whether the learned
policies change their actions in response to the observed state (Section
5.4).

\textbf{O5.} Measure the effect of adding drawdown to the state while
keeping the other experimental conditions fixed (Section 5.5).

\textbf{O6.} Compare the final learned agents with the benchmark
strategies over 24 held-out test months, using six independently trained
random seeds for each configuration (Section 5.6).

\textbf{O7.} Examine whether results change when a different reward
formula is used (Section 5.7).

These objectives create a structured way to evaluate both the financial
performance and the learned behaviour of the three reinforcement
learning methods.

\section{Scope}

The study is deliberately narrowed to one asset to make the experiments
more controlled and easier to analyse.

All experiments use the SPY asset, an exchange-traded fund that tracks
the S\&P 500 index, with daily data from 2018 to 2025. The data are
divided chronologically into training (2018--2022), validation (2023),
and held-out for testing periods (2024--2025).

Each monthly episode starts with \$10,000 of capital and no asset
holdings, and each transaction to open a position incurs a 0.05\% fee.
The available actions are discrete: the agent can buy a fixed fraction
of the initial capital, sell the same amount, or hold. The action
masking prevents invalid actions, such as attempting to sell when no
position is held.

Three reinforcement learning algorithms are evaluated. REINFORCE and
Deep Q-learning are implemented from scratch, while PPO is implemented
using a popular library.

Several areas are intentionally outside the scope of this study. These
include multi-asset portfolios, continuous position sizing, intraday
trading data, and live deployment. Keeping the study within these
boundaries limits the number of factors that can affect the results,
making it easier to identify which specific factor causes a change in
performance or behaviour during testing rather than multiple changes at
once.

\section{Contributions}

This dissertation makes four main contributions:

\textbf{1. A framework for evaluating trading agents.}

The study uses a controlled process in which experiment variables, such
as trade size and transaction cost, are set using validation data and
finalized before moving on to the held-out test data. The trading
environment, state features, rewards, and benchmark strategies are also
verified through tests to identify and correct problems before final
evaluation,

\textbf{2. A behavioural analysis for testing state utilization.}

A separate analysis was developed to determine whether a policy learned
actually responds to the information in its state, rather than simply
following a fixed pattern of available actions.

In this evaluation, steps are grouped by their action mask. When a
policy selects different actions at steps with the same mask pattern,
the difference cannot be explained by the available actions alone,
thereby demonstrating a change in its action based on the current
observed state. The analysis does not identify which individual state
feature drives this dependence, nor does it make claims about the
learned network\textquotesingle s internal representations.

This reveals a difference between algorithms that would not be seen from
financial return figures alone: all six Deep Q-learning seeds were
state-dependent, while none of the twelve policy-gradient seeds
(REINFORCE and PPO) were (Section 5.4).

\textbf{3. An explanation for the negative results.}

None of the three learning algorithms outperformed buy-and-hold across
the 24 held-out test months. The behavioural analysis helps explain this
outcome. The policies that achieved returns closest to the benchmark,
REINFORCE, generally did so through a fixed pattern of high exposure,
while the algorithm that did use its state information, Deep Q-learning,
reduced market participation rather than successfully understanding
market conditions to perform better.

Although the overall result is negative, evidence from the six training
seeds per configuration, behavioural analysis, and peripheral
experiments on state representation, reward function, and transaction
costs supports it as a finding rather than a failure.

\section{Dissertation Structure}

The dissertation is organized as follows. Chapter 2 reviews existing
research on reinforcement learning for financial trading and portfolios
and explains how this study's evaluation approach relates to previous
work.

Chapter 3 presents the mathematical formulation of the portfolio trading
problem, describing the trading environment as a Markov decision
process; the state, action, and reward formulas; the three reinforcement
learning algorithms; and their optimisation functions.

Chapter 4 describes the system implementation, data pipeline,
verification process, trading environment, and the corrections made
during development.

Chapter 5 presents the experiment results and analyses them in relation
to the four research questions. It reports the trade-size calibration,
behavioural analysis, benchmark comparison, state representation and
reward experiment, and transaction-cost calibration.

Finally, Chapter 6 concludes the dissertation by summarising the main
findings, critically evaluating the work, discussing its limitations and
possible directions for future research.
